\documentclass[11pt]{article}
\usepackage{amsmath, amssymb}
\usepackage{amsfonts}
\usepackage{enumitem}
\usepackage{Braket}
\usepackage{geometry}
\usepackage{listings}
\usepackage{graphicx}
\usepackage{physics}
\usepackage{hyperref}
\usepackage{tcolorbox}
\usepackage{bm}
\geometry{a4paper, margin=1in}

\title{My First LaTeX Document}
\author{Your Name}
\date{\today}

\begin{document}
Master Laboratory
Optical Pumping
Group 18
Performed on April 5th-April 9th
under supervision of
Abstract
% Optical pumping is a method of changing occupation numbers of atomic system
% by irradiation of the system with characteristic light. In the experiment this report
% documents the method was used for determining several characteristics of Rubidium
% (85Rb and 87Rb) as well as measuring earth’s magnetic field strength in various
% directions.
% The hyperfine structure of 85Rb and 87Rb was investigated and six out of eigth
% expected spectral lines could be resolved. With the frequency of the hyperfine
% structure energy shifts we were able to calculate three out of four hyperfine structure
% constants A87Rb,2S1/2 
% = (1.44±0.03)×10−5 eV,A87Rb,2P1/2 
% = (1.54±0.27)×10−6 eV,
% and A85Rb,2S1/2 
% = (4.71 ± 0.15) × 10−6eV, the first two being in a 2-σ-environment
% and the third in a 4-σ-environment of the literature value.
% With optical pumping under double resonance with we determined the nuclear
% spins of 85Rb and 87Rb to be I85Rb = 2.51±0.03 and I87Rb = 1.49±0.03, respectively;
% the values are extremely well compatible with the literature values and feature a
% relative error of ≈ 2%. Further, earths magnetic field was measured with optical
% pumping under double resonance to be Bvert = (33.9 ± 0.5)µT,
% Bhor = (5.6 ±
% 1.1) µT which deviates by about 18 standard deviations from the literature value.
% Leveraging the precession of the ensembles spin in the vertical component of
% earth’s magnetic field, the absolute value of the latter was determined to be Bvert =
% (32.6±1.4)µT, deviating 7-σ from the literature value but compatible with the value
% from the double resonance part of the experiment.
% With the methods of Dehmelt and Franzen, the relaxation time of the system
% was determined to be T85Rb
% R =(4.3±2.5)ms,T87Rb
% R =(−5.1±2.4)ms, and TR =
% (8.5 ± 2.5)ms, deviating at most 5-σ from the literature value. While the methods
% could be succesfully replicated, the results are not consistent with each other.
\end{document}